\textbf{Questão 2.9:} Confirme a eq. (2.30) aplicando o restante do teorema Pi de Buckingham ao problema do pêndulo.

\vspace{0,25cm}
\textit{Solução:}
Consideremos as equações iniciais:
\begin{center}
    $\pi_{1} = \frac{T_{0}}{\sqrt{\frac{l}{g}}}$  \\
    \vspace{0,2cm}
    $\pi_{2} = \theta$ \\
     \vspace{0,2cm}
    $\pi_{3} = \frac{T}{m \cdot g}$
\end{center}

Agora, construindo a tabela de dimensões:
\begin{table}[h]
\centering
\begin{tabular}{lcc}
\toprule
\textbf{Grandeza} & \textbf{Símbolo} & \textbf{Dimensão} \\ \midrule
Comprimento & $l$ & $L$  \\
Aceleração gravitacional & $g$ & $LT^{-2}$  \\
Massa & $m$ & $M$  \\
Período & $T_0$ & $T$ \\
Ângulo & $\theta$ & $1$ (adimensional)  \\
Tensão da corda & $T$ & $MLT^{-2}$  \\ \bottomrule
\end{tabular}
\end{table}
    

As dimensões básicas são $L$, $M$, $T$ e temos 6 variáveis. Consideraremos como variáveis básicas $l$, $m$, $g$. 

Além disso, pelo Teorema Pi de Buckingham, teremos 3 grupos $\pi$.

Vamos construir os grupos $\pi$:
\begin{center}
$\pi_{1} = l^{a_1} g^{b_1} m^{c_1} T_{0}$ \\
 \vspace{0,2cm}
$\pi_{2} = l^{a_2} g^{b_2} m^{c_2} \theta$ \\
 \vspace{0,2cm}
$\pi_{3} = l^{a_3} g^{b_3} m^{c_3} T$ \\
\end{center}

Para cada grupo $\pi$ vamos impor a condição que seja adimensional, ou seja, $\pi=M^0L^0T^0$.


Tomando $\pi_{1}$, 
\begin{center}
    $\pi_{1} = l^{a_1} g^{b_1} m^{c_1} T_{0}$ \\
     \vspace{0,2cm}
     $(L)^{a_1} (LT^{-2})^{b_1} (M)^{c_1} (T)^1 = M^0 L^0 T^0$
\end{center}

\begin{itemize}
    \item Para a massa $M$ \\
    $c_{1}=0$
    \item Para o comprimento $L$ \\
    $a_{1}+b_{1}=0 \Leftrightarrow a_{1}=-b_{1}$
    \item Para o tempo $T$ \\
    $-2b_{1}+1=0 \Rightarrow b_{1}=\frac{1}{2}$ \\
Sendo assim, $a_{1}=-\frac{1}{2}$
\end{itemize}

Agora que temos os valores dos expoentes, vamos voltar à fórmula $\pi_1$:

\begin{center}
     $\pi_{1} = l^{a_1} g^{b_1} m^{c_1} T_{0}$ \\
      \vspace{0,2cm}
      $\pi_{1} = l^{-\frac{1}{2}} g^{\frac{1}{2}} m^{0} T_{0}$ \\
       \vspace{0,2cm}
     $\pi_1=(L)^{-\frac{1}{2}} (LT^{-2})^{\frac{1}{2}} (M)^{0} (T)^1$ \\
     \vspace{0,2cm}
     $\pi_1=1$
\end{center}

Com isso, verificamos que a fórmula dada em $\pi_1$ é válida. 

Agora, tomando $\pi_{2}$, 
\begin{center}
    $\pi_{2} = l^{a_2} g^{b_2} m^{c_2}\theta$ \\
     \vspace{0,2cm}
     $(L)^{a_2} (LT^{-2})^{b_2} (M)^{c_2}1 = M^0 L^0 T^0$
\end{center}

\begin{itemize}
    \item Para a massa $M$ \\
    $c_{2}=0$
    \item Para o comprimento $L$ \\
    $a_{2}+b_{2}=0 \Leftrightarrow a_{2}=-b_{2}$
    \item Para o tempo $T$ \\
    $-2b_{2}=0 \Rightarrow b_{2}=0$ \\
Sendo assim, $a_{2}=0$
\end{itemize}

Agora que temos os valores dos expoentes, vamos voltar à fórmula $\pi_2$:

\begin{center}
     $\pi_{2} = l^{a_2} g^{b_2} m^{c_2} \theta$ \\
      \vspace{0,2cm}
      $\pi_{2} = l^{0} g^{0} m^{0}\theta$ \\
       \vspace{0,2cm}
     $\pi_{2}=(L)^{0} (LT^{-2})^{0} (M)^{0} 1$ \\
     \vspace{0,2cm}
     $\pi_{2}=1$
\end{center}

Com isso, verificamos que a fórmula dada em $\pi_{2}$ é válida.

Por fim, tomando $\pi_{3}$, 
\begin{center}
    $\pi_{3} = l^{a_1} g^{b_1} m^{c_1} T$ \\
     \vspace{0,2cm}
     $(L)^{a_3} (LT^{-2})^{b_3} (M)^{c_3} (MLT^{-2}) = M^0 L^0 T^0$
\end{center}

\begin{itemize}
    \item Para a massa $M$ \\
    $c_{3}+1=0 \Rightarrow c_{3}=-1$
    \item Para o comprimento $L$ \\
    $a_{3}+b_{3}+1=0 \Leftrightarrow a_{3}+b_{3}=-1$
    \item Para o tempo $T$ \\
    $-2b_{3}-2=0 \Rightarrow b_{3}=-1$ \\
Sendo assim, \\
$a_{3}+b_{3}=-1 \Rightarrow a_{3}-1=-1 \Rightarrow a_{3}=0$
\end{itemize}

Agora que temos os valores dos expoentes, vamos voltar à fórmula $\pi_{3}$:

\begin{center}
     $\pi_{3} = l^{a_3} g^{b_3} m^{c_3} T$ \\
      \vspace{0,2cm}
      $\pi_{3} = l^{0} g^{-1} m^{-1} T$ \\
       \vspace{0,2cm}
     $\pi_{3}=(L)^{0} (LT^{-2})^{-1} (M)^{-1} (MLT^{-2})$ \\
     \vspace{0,2cm}
     $\pi_{3}=1$
\end{center}

Com isso, verificamos que a fórmula dada em $\pi_3$ é válida.